Proyek \textit{Capstone SCADA-Based Optimal Load Shedding and Recovery} ini mendukung keandalan dan stabilitas sistem tenaga listrik sub-transmisi di tengah tingginya penetrasi Energi Baru Terbarukan (EBT) yang berdampak pada penurunan inersia sistem. Sesuai dengan standar operasi jaringan (\textit{Grid Code}) seperti Peraturan Menteri ESDM No. 3 Tahun 2007 dan IEEE Std 1547-2018, frekuensi sistem wajib dijaga dalam rentang aman untuk mencegah pemadaman total (\textit{blackout}). Salah satu mekanisme pertahanan darurat yang umum digunakan saat ini adalah skema \textit{Under Frequency Load Shedding} (UFLS) konvensional, namun metode statis ini rentan memicu \textit{overshedding} (pelepasan beban berlebih). Sistem ini dirancang untuk mengatasi kelemahan tersebut dengan mendeteksi gangguan secara \textit{real-time} melalui pembacaan telemetri deviasi frekuensi dan \textit{Rate of Change of Frequency} (RoCoF) yang dimodelkan menggunakan \textit{physics engine} berbasis \textit{swing equation}. Keputusan pemutusan beban dihitung menggunakan metode optimasi matematis \textit{Mixed-Integer Linear Programming} (MILP-DC) yang menentukan jumlah beban minimal yang harus dilepas untuk menstabilkan sistem. Nilai yang menjadi referensi optimasi memastikan bahwa kerugian ekonomi (\textit{Cost of Energy Not Served}/CENS) ditekan seminimal mungkin dengan tetap menegakkan batasan fisik (\textit{hard constraint}) seperti profil tegangan AC, kapasitas saluran, dan perlindungan absolut terhadap beban kritis. Sistem ini divalidasi melalui pendekatan \textit{Hardware-in-the-Loop} (HIL) yang memanfaatkan \textit{Programmable Logic Controller} (PLC) Schneider Electric Modicon M221 sebagai \textit{Remote Terminal Unit} (RTU). Akses komunikasi dijalankan secara \textit{real-time} menggunakan protokol standar industri Modbus TCP/IP antara perangkat keras dan algoritma \textit{backend} Python. Hasil deteksi optimasi akan secara otomatis mengeksekusi \textit{circuit breaker} (CB) beban untuk fase pemutusan (\textit{shedding}) maupun penjadwalan penyalaan beban kembali secara bertahap (\textit{recovery}). Seluruh proses otomasi ini dicatat dalam log \textit{Sequence of Events} (SOE) untuk keperluan audit. Jika sistem deteksi tidak berjalan sebagaimana mestinya atau kondisi mengharuskan intervensi, operator (\textit{dispatcher}) dapat memantau status jaringan melalui \textit{Single Line Diagram} (SLD) dinamis dan melakukan \textit{override} manual melalui \textit{Web HMI Dashboard}.